\documentclass[11pt]{article}
\usepackage[a4paper,margin=1in]{geometry}
\usepackage{amsmath,amssymb,amsfonts}
\usepackage{booktabs}
\usepackage{graphicx}
\usepackage{caption}
\usepackage{enumitem}
\usepackage[table]{xcolor}
\usepackage{array}
\usepackage[hidelinks]{hyperref}
\usepackage{xurl}
\usepackage{microtype}

\title{2D Corridor Membranes and Near-Start Two-Target Bimodality\\\large New Extension Report Bridging Previous 2D Results}
\author{}
\date{}

\setlist{nosep}
\setlength{\textfloatsep}{8pt plus 2pt minus 2pt}
\setlength{\floatsep}{6pt plus 2pt minus 2pt}
\setlength{\intextsep}{6pt plus 2pt minus 2pt}
\setlength{\abovecaptionskip}{4pt}
\setlength{\belowcaptionskip}{0pt}
\setlength{\tabcolsep}{5pt}
\renewcommand{\arraystretch}{1.10}

\begin{document}
\maketitle

\section*{Headline Results}
This report adds two new research tracks on the 2D rectangular setting while keeping the same phase criteria and peak-valley style used in earlier reports.
\begin{enumerate}
  \item \textbf{Corridor + one target:} replace fully reflecting corridor walls by semi-permeable membranes, with both symmetric ($\kappa_\uparrow=\kappa_\downarrow$) and asymmetric ($\kappa_\uparrow\neq\kappa_\downarrow$) cases.
  \item \textbf{No corridor + two targets:} place one target close to the start and keep the other at the far side, to test robust ``fast near-hit + slow far-hit'' splitting.
\end{enumerate}

Main outcomes:
\begin{itemize}
  \item For symmetric membranes, clear-double forms a continuous ultra-low leakage band (see Figure~\ref{fig:sym-phase} and Table~\ref{tab:sym-w16}); it then enters a weak-double transition band and finally turns mostly single-peak at higher leakage.
  \item A clear asymmetric representative ($\kappa_\uparrow=0.002,\kappa_\downarrow=0$) keeps robust bimodality and changes valley/second-peak composition in the L$\times$R class decomposition.
  \item In the no-corridor two-target setup, near-start placement yields stable clear-double: the representative case $(b_x,d,\Delta y)=(0.07,2,3)$ gives near/far mass $\approx0.326/0.674$ with peaks around $(11,739)$; the dedicated showcase clear-instance is $(0.12,4,1)$ with near/far mass $\approx0.350/0.650$ and a deeper valley (valley/max $\approx0.007$).
\end{itemize}

\section{Model and Shared Baseline}
We use the existing rectangular 2D framework with reflecting boundaries:
\[
L_x=60,\quad q=0.8,\quad x_s=7,\quad x_t=58.
\]
For the corridor branch we use an updated anchor tuned for higher clear-double coverage:
\[
h=2,\quad m=5,\quad \delta_{\rm core}=1.00,\quad \delta_{\rm open}=0.55,\quad b_x=-0.08.
\]

\subsection{Membrane Definition}
Each wall-crossing attempt at a membrane edge is accepted with probability $\kappa$ and reflected with probability $1-\kappa$.
Symmetric case: $\kappa_\uparrow=\kappa_\downarrow=\kappa$.
Asymmetric case: $\kappa_\uparrow\neq\kappa_\downarrow$.

\subsection{Class Labels for Mechanism}
We keep peak/valley windows and use a compact 4-class coding:
\begin{itemize}
  \item $L$: membrane leak happened or not;
  \item $R$: committor recross happened or not.
\end{itemize}
Classes are \texttt{L0R0}, \texttt{L0R1}, \texttt{L1R0}, \texttt{L1R1}.

\section{Track A: Corridor + One Target + Membranes}
\subsection{Symmetric Scan}
\begin{figure}[!htbp]
  \centering
  \includegraphics[width=0.86\linewidth]{figures/membrane_symmetric_phase_map.pdf}
  \caption{Symmetric membrane phase map (rows: $W_y$, columns: $\kappa$).}
  \label{fig:sym-phase}
\end{figure}

\begin{figure}[!htbp]
  \centering
  \includegraphics[width=0.86\linewidth]{figures/membrane_symmetric_sep_map.pdf}
  \caption{Symmetric membrane separation map (sep-score).}
  \label{fig:sym-sep}
\end{figure}

\begin{table}[!htbp]
  \centering
  \footnotesize
  \begin{tabular}{@{}cccccc@{}}
\toprule
$\kappa$ & phase & $t_{p1}$ & $t_{p2}$ & valley/max & sep-score \\
\midrule
0 & 2 & 395 & 1711 & 0.058 & 1.87 \\
0.0003 & 2 & 395 & 1710 & 0.061 & 1.84 \\
0.0005 & 2 & 395 & 1708 & 0.063 & 1.80 \\
0.0008 & 2 & 395 & 1706 & 0.066 & 1.76 \\
0.001 & 2 & 394 & 1704 & 0.069 & 1.72 \\
0.0013 & 2 & 394 & 1702 & 0.072 & 1.68 \\
0.0015 & 2 & 394 & 1700 & 0.074 & 1.64 \\
0.0018 & 2 & 394 & 1698 & 0.077 & 1.59 \\
0.002 & 2 & 394 & 1696 & 0.080 & 1.53 \\
0.0022 & 2 & 393 & 1694 & 0.082 & 1.46 \\
0.0025 & 2 & 393 & 1692 & 0.085 & 1.40 \\
0.0027 & 2 & 393 & 1690 & 0.087 & 1.34 \\
0.003 & 2 & 393 & 1689 & 0.090 & 1.29 \\
0.0032 & 2 & 393 & 1687 & 0.093 & 1.25 \\
0.0035 & 2 & 393 & 1685 & 0.095 & 1.21 \\
0.0037 & 1 & 392 & 1683 & 0.098 & 0.84 \\
0.004 & 1 & 392 & 1681 & 0.100 & 0.84 \\
0.0045 & 1 & 392 & 1677 & 0.105 & 0.82 \\
0.005 & 1 & 392 & 1673 & 0.110 & 0.81 \\
0.006 & 1 & 391 & 1664 & 0.120 & 0.79 \\
0.0075 & 1 & 390 & 1652 & 0.134 & 0.77 \\
0.01 & 1 & 388 & 1629 & 0.157 & 0.73 \\
0.015 & 1 & 386 & 1576 & 0.198 & 0.67 \\
0.02 & 1 & 383 & 1505 & 0.235 & 0.61 \\
0.03 & 0 & - & - & - & 0.00 \\
0.05 & 0 & - & - & - & 0.00 \\
0.1 & 0 & - & - & - & 0.00 \\
0.2 & 0 & - & - & - & 0.00 \\
\bottomrule
\end{tabular}

  \caption{Symmetric membrane metrics on the $W_y=16$ slice.}
  \label{tab:sym-w16}
\end{table}

In this run (taking the $W_y=16$ slice as reference), strong bimodality gradually weakens with permeability: clear-double for $\kappa\in[0,0.0035]$, weak-double around $\kappa\in[0.0037,0.02]$, and mostly single-peak for $\kappa\ge 0.03$.

\subsection{Asymmetric Scan}
\begin{figure}[!htbp]
  \centering
  \includegraphics[width=0.86\linewidth]{figures/membrane_asymmetric_phase_map.pdf}
  \caption{Asymmetric membrane phase map ($\kappa_\uparrow$ vs $\kappa_\downarrow$).}
  \label{fig:asym-phase}
\end{figure}

\begin{figure}[!htbp]
  \centering
  \includegraphics[width=0.86\linewidth]{figures/membrane_asymmetric_sep_map.pdf}
  \caption{Asymmetric membrane separation map.}
  \label{fig:asym-sep}
\end{figure}

\begin{table}[!htbp]
  \centering
  \footnotesize
  \begin{tabular}{@{}cccccc@{}}
\toprule
$\kappa_{\uparrow}$ & $\kappa_{\downarrow}$ & phase & $t_{p1}$ & $t_{p2}$ & sep-score \\
\midrule
0 & 0.0003 & 2 & 395 & 1711 & 1.85 \\
0.0003 & 0 & 2 & 395 & 1710 & 1.85 \\
0.0005 & 0 & 2 & 395 & 1710 & 1.84 \\
0 & 0.0005 & 2 & 395 & 1710 & 1.84 \\
0.0008 & 0 & 2 & 395 & 1709 & 1.82 \\
0.0005 & 0.0003 & 2 & 395 & 1709 & 1.82 \\
0.0003 & 0.0005 & 2 & 395 & 1709 & 1.82 \\
0 & 0.0008 & 2 & 395 & 1709 & 1.82 \\
0.001 & 0 & 2 & 395 & 1708 & 1.80 \\
0.002 & 0 & 2 & 394 & 1704 & 1.73 \\
\bottomrule
\end{tabular}

  \caption{Top asymmetric membrane cases (ranked by phase and separation).}
  \label{tab:asym-top}
\end{table}

Representative overlays:
\begin{figure}[!htbp]
  \centering
  \includegraphics[width=0.88\linewidth]{figures/membrane_rep_sym_geometry.pdf}
  \caption{Symmetric membrane representative geometry (corridor band, membrane location, start/target, local-bias direction).}
  \label{fig:rep-sym-geo}
\end{figure}

\begin{figure}[!htbp]
  \centering
  \includegraphics[width=0.86\linewidth]{figures/membrane_rep_sym_fpt.pdf}
  \caption{Symmetric representative ($\kappa=0.002$): total FPT and corridor/outer-source flux decomposition.}
  \label{fig:rep-sym-fpt}
\end{figure}

\begin{figure}[!htbp]
  \centering
  \includegraphics[width=0.88\linewidth]{figures/membrane_rep_asym_geometry.pdf}
  \caption{Asymmetric membrane representative geometry (different top/bottom membrane permeabilities).}
  \label{fig:rep-asym-geo}
\end{figure}

\begin{figure}[!htbp]
  \centering
  \includegraphics[width=0.86\linewidth]{figures/membrane_rep_asym_fpt.pdf}
  \caption{Asymmetric representative ($\kappa_\uparrow=0.002,\kappa_\downarrow=0$): total FPT and flux decomposition.}
  \label{fig:rep-asym-fpt}
\end{figure}

\begin{figure}[!htbp]
  \centering
  \includegraphics[width=0.94\linewidth]{figures/membrane_class_window_bars.pdf}
  \caption{Window-wise composition of 4 classes (L$\times$R).}
  \label{fig:membrane-class-bars}
\end{figure}

\begin{figure}[!htbp]
  \centering
  \includegraphics[width=0.90\linewidth]{figures/membrane_class_legend.pdf}
  \caption{Representative L$\times$R legend: left, 2x2 class definition; right, typical dominant-class transfer across peak1/valley/peak2 windows.}
  \label{fig:membrane-class-legend}
\end{figure}

\begin{table}[!htbp]
  \centering
  \footnotesize
  \begin{tabular}{@{}lccccc@{}}
\toprule
Window & Case & L0R0 & L0R1 & L1R0 & L1R1 \\
\midrule
peak1 & sym & 0.996 & 0.000 & 0.004 & 0.000 \\
peak1 & dir & 0.998 & 0.000 & 0.002 & 0.000 \\
valley & sym & 0.581 & 0.116 & 0.299 & 0.004 \\
valley & dir & 0.655 & 0.121 & 0.224 & 0.001 \\
peak2 & sym & 0.142 & 0.742 & 0.080 & 0.036 \\
peak2 & dir & 0.166 & 0.749 & 0.067 & 0.018 \\
\bottomrule
\end{tabular}

  \caption{Numerical values of the 4-class window composition.}
  \label{tab:membrane-window-classes}
\end{table}

\begin{table}[!htbp]
  \centering
  \footnotesize
  \begin{tabular}{@{}lllc c@{}}
\toprule
Case & $q^*$ & Window & dominant class & $w_{\mathrm{L1R1}}$ \\
\midrule
sym & 0.4 & peak1 & L0R0 & 0.000 \\
sym & 0.4 & valley & L0R0 & 0.197 \\
sym & 0.4 & peak2 & L0R1 & 0.104 \\
dir & 0.4 & peak1 & L0R0 & 0.000 \\
dir & 0.4 & valley & L0R0 & 0.196 \\
dir & 0.4 & peak2 & L0R1 & 0.082 \\
sym & 0.5 & peak1 & L0R0 & 0.000 \\
sym & 0.5 & valley & L0R0 & 0.004 \\
sym & 0.5 & peak2 & L0R1 & 0.036 \\
dir & 0.5 & peak1 & L0R0 & 0.000 \\
dir & 0.5 & valley & L0R0 & 0.001 \\
dir & 0.5 & peak2 & L0R1 & 0.018 \\
sym & 0.6 & peak1 & L0R0 & 0.000 \\
sym & 0.6 & valley & L0R0 & 0.002 \\
sym & 0.6 & peak2 & L0R1 & 0.022 \\
dir & 0.6 & peak1 & L0R0 & 0.000 \\
dir & 0.6 & valley & L0R0 & 0.000 \\
dir & 0.6 & peak2 & L0R0 & 0.008 \\
\bottomrule
\end{tabular}

  \caption{$q^*$ robustness check (exact decomposition): dominant class and $w_{\mathrm{L1R1}}$ across thresholds.}
  \label{tab:membrane-qstar-sens}
\end{table}

With strict recross (visit $\Sigma_{q^*}$ then return to $A$), the exact decomposition shows: peak1 is \texttt{L0R0}-dominated, valley remains mostly \texttt{L0R0} with increased \texttt{L1R0}, and peak2 is \texttt{L0R1}-dominated. So the late peak is mainly a commitment-failure effect, not necessarily leak-driven.

\section{Track B: No Corridor + Two Targets (Near-Start Target)}
We remove corridor geometry and keep reflecting boundaries only. One target is placed near the start, and we optimize in two stages: first maximize clear coverage over candidate grids, then select a dedicated clear-instance from all candidate results for the sharpest double-peak visualization.

\begin{figure}[!htbp]
  \centering
  \includegraphics[width=0.88\linewidth]{figures/two_target_nearstart_phase_map.pdf}
  \caption{No-corridor two-target phase map (near-target distance $d$ vs $b_x$).}
  \label{fig:tt-near-phase}
\end{figure}

\begin{figure}[!htbp]
  \centering
  \includegraphics[width=0.88\linewidth]{figures/two_target_nearstart_sep_map.pdf}
  \caption{No-corridor two-target separation map.}
  \label{fig:tt-near-sep}
\end{figure}

\begin{table}[!htbp]
  \centering
  \footnotesize
  \begin{tabular}{@{}cccccc@{}}
\toprule
$b_x$ & clear $d$ interval & best $d$ & phase(best) & $P(\mathrm{near})$ & $P(\mathrm{far})$ \\
\midrule
+0.00 & - & 3 & 0 & 0.892 & 0.108 \\
+0.02 & - & 2 & 0 & 0.603 & 0.397 \\
+0.04 & 3-9 & 3 & 2 & 0.444 & 0.556 \\
+0.05 & 2-10 & 2 & 2 & 0.393 & 0.607 \\
+0.06 & 2-10 & 2 & 2 & 0.355 & 0.645 \\
+0.07 & 2-10 & 2 & 2 & 0.326 & 0.674 \\
+0.08 & 2-10 & 2 & 2 & 0.301 & 0.699 \\
+0.09 & 2-10 & 2 & 2 & 0.281 & 0.719 \\
+0.10 & 2-10 & 2 & 2 & 0.264 & 0.736 \\
+0.11 & 2-10 & 2 & 2 & 0.248 & 0.752 \\
+0.12 & 2-10 & 2 & 2 & 0.235 & 0.765 \\
+0.14 & 2-10 & 2 & 2 & 0.212 & 0.788 \\
+0.16 & 2-10 & 3 & 2 & 0.206 & 0.794 \\
\bottomrule
\end{tabular}

  \caption{Per-$b_x$ summary: clear intervals and best near-target distance.}
  \label{tab:tt-near-summary}
\end{table}

Representative case ($W_y=16$, $b_x=0.07$, $d=2$):
\begin{figure}[!htbp]
  \centering
  \includegraphics[width=0.88\linewidth]{figures/two_target_rep_geometry.pdf}
  \caption{Representative no-corridor geometry (start, near target, far target, and global drift direction).}
  \label{fig:tt-near-rep-geo}
\end{figure}

\begin{figure}[!htbp]
  \centering
  \includegraphics[width=0.86\linewidth]{figures/two_target_rep_fpt.pdf}
  \caption{Representative no-corridor case: total FPT with near/far splitting components.}
  \label{fig:tt-near-rep-fpt}
\end{figure}

\begin{figure}[!htbp]
  \centering
  \includegraphics[width=0.94\linewidth]{figures/two_target_rep_doublepeak_example.pdf}
  \caption{No-corridor clear-double exemplar: left, a tail-emphasized linear view (peak1 is clipped and marked as \texttt{p1*} for readability of valley/peak2); right, a log view with $y\ge 10^{-6}$.}
  \label{fig:tt-near-doublepeak-example}
\end{figure}

\begin{table}[!htbp]
  \centering
  \footnotesize
  \resizebox{\linewidth}{!}{\begin{tabular}{@{}cccccccccccccccc@{}}
\toprule
spec & $b_x$ & $d$ & $\Delta y$ & phase & sep-score & valley/max & $p_{\min}$ & peak-ratio & peak-margin & min-margin & gate-pass & gate-slack & clear-score & showcase & $\varepsilon_{\mathrm{split}}$ \\
\midrule
adaptive\_refine & +0.12 & 4 & 1 & 2 & 2.38 & 0.007 & 0.350 & 0.170 & 0.070 & 0.070 & yes & 0.070 & 1.54 & 1.92 & $1.18\times 10^{-9}$ \\
\bottomrule
\end{tabular}
}
  \caption{No-corridor clear-instance quality metrics (auto-selected from all candidate scans). We report \texttt{peak-margin}, \texttt{min-margin}, and strong-gate audit fields (\texttt{gate-pass}/\texttt{gate-slack}) to avoid near-threshold pseudo-clear picks, and include \texttt{showcase} for visual ranking consistency.}
  \label{tab:tt-near-clear-instance}
\end{table}

\begin{figure}[!htbp]
  \centering
  \includegraphics[width=0.94\linewidth]{figures/two_target_no_corridor_clear_instance.pdf}
  \caption{Dedicated no-corridor clear-instance figure (maximum showcase-score over all candidate scans): tail-emphasized linear panel plus log panel to keep valley/peak2 structure explicit.}
  \label{fig:tt-near-clear-instance}
\end{figure}

\begin{figure}[!htbp]
  \centering
  \includegraphics[width=0.82\linewidth]{figures/two_target_rep_window_split.pdf}
  \caption{Window-wise near/far mass fractions in the representative case.}
  \label{fig:tt-near-window}
\end{figure}

The clear-instance metrics (near/far masses, sep-score, valley/max, peak/min margins, showcase score, and splitting-committor closure error $\varepsilon_{\rm split}$) are listed in Table~\ref{tab:tt-near-clear-instance}. In this round we apply a robustness gate before final ranking (prioritizing positive \texttt{peak-margin/min-margin} with deep-valley preference), so the selected no-corridor case is not only phase=2-valid but also visually stable under threshold perturbations. The current auto-selected clear-instance comes from the adaptive-refinement round, $(b_x,d,\Delta y)=(0.12,4,1)$, with \texttt{peak-margin}=0.070, \texttt{min-margin}=0.070, and \texttt{showcase}=1.92. It passes the first gate:
\[
\mathrm{peak\mbox{-}margin}\ge 0.07,\qquad \mathrm{min\mbox{-}margin}\ge 0.07,\qquad \mathrm{valley}/\max\le 0.10.
\]
Its primary-gate audit reports \texttt{gate-pass=yes}, \texttt{gate-slack}=0.070 (minimum slack over the three thresholds), and 168 candidates in the first gate pool.

\section{Method Reference (Committor + 2x2 Classes)}
\subsection*{Reproducible Definition of Phase and sep-score}
To make phase maps fully reproducible, we use the same criteria as in \path{src/vkcore/grid2d/rect_bimodality/cli.py}:
\paragraph{(1) Two-peak detection (\texttt{find\_two\_peaks})}
Let $\tilde f=\mathrm{smooth}_7(f)$ (window-7 smoothing). Candidate peaks satisfy
\[
\tilde f(t)>\tilde f(t-1),\quad \tilde f(t)\ge \tilde f(t+1),\quad \tilde f(t)\ge 0.06\,\max \tilde f.
\]
Pair selection then enforces (default code thresholds)
\[
t_{p2}-t_{p1}\ge 20,\ \mathrm{peak\_balance}\ge 0.08,\ \mathrm{valley\_drop}\ge 0.08,\ \mathrm{minor\_drop}\ge 0.005,
\]
plus side-dominance tolerance (\texttt{dominance\_tol}=0.02). Define
\[
\mathrm{has\_double}:=\mathbf 1\{t_{p1},t_{p2}\ \text{both exist}\}.
\]

\paragraph{(2) One-target phase definition}
Define
\[
\mathrm{sep}_{\rm peaks}:=\frac{|t_{p2}-t_{p1}|}{w_{1/2}^{(1)}+w_{1/2}^{(2)}},\quad
\mathrm{valley}/\max:=\frac{h_v}{\max(h_{p1},h_{p2})},\quad
\mathrm{peak\_balance}:=\frac{\min(h_{p1},h_{p2})}{\max(h_{p1},h_{p2})}.
\]
Here $w_{1/2}^{(k)}$ is half-width at half-maximum of peak $k$ (computed on raw $f$), while $h_{p1},h_{p2},h_v$ are read from smoothed $\tilde f$.
Then
\[
\mathrm{phase}_{\rm OT}=
\begin{cases}
2,& \mathrm{has\_double}=1,\ \mathrm{sep}_{\rm peaks}\ge 1.0,\ \mathrm{valley}/\max\le 0.35,\ \mathrm{peak\_balance}\ge 0.15,\\
1,& \mathrm{has\_double}=1\ \text{but the above clear-double test fails},\\
0,& \mathrm{has\_double}=0.
\end{cases}
\]

\paragraph{(3) Two-target phase definition}
Define modal times and mode-separation score
\[
t_{\rm mode}^{\rm near}:=\arg\max_t f_{\rm near}(t),\quad
t_{\rm mode}^{\rm far}:=\arg\max_t f_{\rm far}(t),
\]
\[
\mathrm{sep}_{\rm mode}:=\frac{|t_{\rm mode}^{\rm far}-t_{\rm mode}^{\rm near}|}{w_{1/2}^{\rm far}+w_{1/2}^{\rm near}}.
\]
Define normalized splitting masses
\[
P_{\rm near}:=\frac{\sum_t f_{\rm near}(t)}{\sum_t f_{\rm any}(t)},\qquad
P_{\rm far}:=\frac{\sum_t f_{\rm far}(t)}{\sum_t f_{\rm any}(t)}.
\]
On the total law, define
\[
\mathrm{peak\_ratio}:=\frac{\min(h_{p1},h_{p2})}{\max(h_{p1},h_{p2})},\qquad
\mathrm{valley}/\max:=\frac{h_v}{\max(h_{p1},h_{p2})}.
\]
Then
\[
\mathrm{phase}_{\rm TT}=
\begin{cases}
2,& \mathrm{has\_double}=1,\ \mathrm{sep}_{\rm mode}\ge 1.0,\ \min(P_{\rm near},P_{\rm far})\ge 0.15,\\
& \mathrm{valley}/\max\le 0.35,\ \mathrm{peak\_ratio}\ge 0.10,\ (\mathrm{peak\_ratio}-\mathrm{valley}/\max)\ge 0.01,\\
1,& \mathrm{has\_double}=1\ \text{but the above clear-double test fails},\\
0,& \mathrm{has\_double}=0.
\end{cases}
\]
So mathematically, \textbf{phase=1} means ``double peaks are detected, but at least one clear-double threshold is not met''.

\subsection*{Committor Mechanism (Detailed)}
Let the hitting time be $\tau_C:=\inf\{t\ge 0:\,X_t\in C\}$. On the discrete Markov chain with transition matrix $P$, let $A$ be the start basin (here a strip near the start) and $B$ be the absorbing target set. Define
\[
q_i := \mathbb P_i(\tau_B < \tau_A),
\]
which solves the discrete Dirichlet boundary-value problem
\[
q_i=0\ (i\in A),\qquad q_i=1\ (i\in B),\qquad
q_i=\sum_j P_{ij} q_j\ (i\notin A\cup B).
\]
With $U=\Omega\setminus(A\cup B)$, this is equivalent to the linear system
\[
(I-P_{UU})\,q_U=P_{UB}\,\mathbf 1.
\]
For this absorbing setup the solution is unique, so $q$ is a dynamics-determined progress coordinate (not an ad hoc geometric proxy).
We use $q$ as the progress coordinate and set $q^*=0.5$ to define the commitment surface
\[
\Sigma_{q^*}:=\{i:\ q_i\ge q^*\}.
\]
Moreover, when $\rho(P_{UU})<1$, the Neumann expansion gives
\[
q_U=(I-P_{UU})^{-1}P_{UB}\mathbf 1
=\sum_{n\ge 0}P_{UU}^n P_{UB}\mathbf 1.
\]
So $q_U$ is the cumulative contribution of paths that spend $n$ interior steps before absorption into $B$, which provides a direct linear-system bridge to the renewal failure-retry interpretation used later: early peak mass is small-$n$ dominated, while the slow peak accumulates from larger-$n$ loops.

\subsection*{L$\times$R Events and Recross Counts}
For a trajectory $X_t$ up to $\tau_B$, define membrane-cross count
\[
N_{\rm mem}:=\sum_{t=0}^{\tau_B-1}\mathbf 1\!\{(X_t,X_{t+1})\in\mathcal E_{\rm mem}\},
\]
and commitment-to-start recross count
\[
N_{\rm ret}:=\sum_{t=0}^{\tau_B-1}\mathbf 1\!\{X_t\in\Sigma_{q^*},\,X_{t+1}\in A\}.
\]
Then use binary labels
\[
L=\mathbf 1\{N_{\rm mem}\ge 1\},\qquad R=\mathbf 1\{N_{\rm ret}\ge 1\},
\]
which yields four classes
\[
(\!L,R\!)\in\{(0,0),(0,1),(1,0),(1,1)\},
\]
implemented as \texttt{L0R0/L0R1/L1R0/L1R1} (Figure~\ref{fig:membrane-class-legend}).

\subsection*{Peak/Valley Mechanism as Renewal Decomposition}
Let $f(t)$ be the FPT law. With direct-success kernel $g(t)$ and fail-return kernel $h(t)$,
\[
f(t)=g(t)+(h*f)(t)=\sum_{k\ge 0}(h^{*k}*g)(t).
\]
Here $k$ is the number of failed attempts before success. Bimodality appears when the $k=0$ (near one-shot) and $k\ge1$ parts separate in both timescale and mass. Window fractions
\[
w_{\ell r}(W)=\frac{\sum_{t\in W} f_{\ell r}(t)}{\sum_{t\in W} f(t)}
\]
make the transfer explicit: peak1 is mainly \texttt{L0R0}; valley remains \texttt{L0R0}-major with a clear rise of \texttt{L1R0}; and peak2 is \texttt{L0R1}-dominant, with \texttt{L1R1} secondary (Figure~\ref{fig:membrane-class-bars}, Table~\ref{tab:membrane-window-classes}).

\subsection*{Closed-Form Bridge from Linear System to FPT Kernels}
For the no-corridor two-target branch, let $U$ denote interior transient states, $e_s$ the interior basis vector at the start, and
\[
Q:=P_{UU},\qquad R_{\rm near}:=P_{U,\partial_{\rm near}},\qquad R_{\rm far}:=P_{U,\partial_{\rm far}}.
\]
Then the first-hit pmfs are
\[
f_{\rm near}(t)=e_s^\top Q^{t-1}R_{\rm near}\mathbf 1,\qquad
f_{\rm far}(t)=e_s^\top Q^{t-1}R_{\rm far}\mathbf 1,\qquad t\ge 1.
\]
Hence the far-hit mass is
\[
P_{\rm far}=\sum_{t\ge 1}f_{\rm far}(t)=e_s^\top(I-Q)^{-1}R_{\rm far}\mathbf 1.
\]
This gives the exact bridge between time-domain decomposition and the linear-system view used by committors.

\subsection*{Conditioned Kernel (Doob Transform) and Peak Transfer}
On interior states with $q_{\rm far}(i)>0$, define
\[
\widetilde P^{\rm far}_{ij}:=\frac{P_{ij}\,q_{\rm far}(j)}{q_{\rm far}(i)}.
\]
This is the Doob $h$-transform of $P$ by $h=q_{\rm far}$, i.e., the exact dynamics conditioned on eventually hitting \emph{far} before \emph{near}. In this view, the near branch is the complementary conditioned dynamics through $q_{\rm near}=1-q_{\rm far}$. Therefore, the window transfer from peak1 to peak2 can be interpreted as a measure transfer between two conditioned ensembles: early-time mass is dominated by short successful excursions, while late-time mass is amplified by loops that repeatedly leave and re-enter low-$q_{\rm far}$ regions before final absorption. This provides a direct dynamical meaning for the observed renewal split.

\subsection*{Spectral View: Origin of Fast/Slow Timescales}
On the interior kernel $Q$, write the (generalized) spectral expansion
\[
Q^{t-1}=\sum_k \lambda_k^{\,t-1} r_k\ell_k^\top,\qquad |\lambda_1|>|\lambda_2|\ge\cdots.
\]
Then
\[
f_{\rm near}(t)=\sum_k c_k^{\rm near}\lambda_k^{\,t-1},\qquad
f_{\rm far}(t)=\sum_k c_k^{\rm far}\lambda_k^{\,t-1}.
\]
Modes with smaller $|\lambda_k|$ decay quickly and dominate early windows, while the leading mode $\lambda_1$ controls the slow tail and late peak. If near/far branches project differently on these modes, we naturally obtain ``near-dominant early window'' versus ``far-dominant late window'' transfer. This is fully consistent with the renewal picture: peak1 is short-stay/low-loop mass, while peak2 is quasi-stationary release from the leading mode. Equivalently, the conditioned hazard
\[
h_{\rm far}(t):=\frac{f_{\rm far}(t)}{S(t-1)}
\]
increases over time when the far channel progressively takes over late windows.

\subsection*{Splitting Committor Closure for Two Targets}
On the same system, define
\[
q_{\rm far}(i):=\mathbb P_i(\tau_{\rm far}<\tau_{\rm near}),\qquad q_{\rm near}(i)=1-q_{\rm far}(i).
\]
Its boundary-value problem is
\[
q_{\rm far}=0\ \text{on near},\qquad q_{\rm far}=1\ \text{on far},\qquad
q_{\rm far}(i)=\sum_j P_{ij}\,q_{\rm far}(j)\ \text{on interior}.
\]
This gives a direct closure with the exact decomposition $f=f_{\rm near}+f_{\rm far}$: $q_{\rm far}(x_s)$ should match $P_{\rm far}$. Define
\[
\varepsilon_{\rm split}:=\left|q_{\rm far}(x_s)-\sum_{t\ge 1}f_{\rm far}(t)\right|
=\left|q_{\rm far}(x_s)-P_{\rm far}\right|.
\]
Moreover, let $\tau:=\tau_{\rm near}\wedge\tau_{\rm far}$ and define
\[
M_t:=q_{\rm far}(X_{t\wedge\tau}).
\]
Because $q_{\rm far}$ is harmonic on interior states and fixed to $0/1$ on near/far boundaries, $M_t$ is a bounded martingale. Optional stopping then gives
\[
q_{\rm far}(x_s)=\mathbb E[M_\tau]
=\mathbb P_{x_s}(\tau_{\rm far}<\tau_{\rm near})
=P_{\rm far}.
\]
So $\varepsilon_{\rm split}$ is theoretically zero; numerically, its scale should track the Bellman-residual scale, making it a joint diagnostic for mechanism closure and solver stability.
\begin{table}[!htbp]
  \centering
  \footnotesize
  \resizebox{\linewidth}{!}{\begin{tabular}{@{}lcccccc@{}}
\toprule
case & $q_{\mathrm{far}}(x_s)$ & $P_{\mathrm{far}}$ & $\varepsilon_{\mathrm{split}}$ & $\|r_{\mathrm{int}}\|_\infty$ & $\|r_{\mathrm{near}}\|_\infty$ & $\|r_{\mathrm{far}}\|_\infty$ \\
\midrule
representative & 0.676009 & 0.676009 & $1.14\times 10^{-9}$ & $9.93\times 10^{-12}$ & $0$ & $0$ \\
clear-instance & 0.650264 & 0.650264 & $1.18\times 10^{-9}$ & $9.93\times 10^{-12}$ & $0$ & $0$ \\
\bottomrule
\end{tabular}
}
  \caption{Splitting-committor closure QC for no-corridor two-target runs (representative + dedicated clear-instance): we report $q_{\rm far}(x_s)$, $P_{\rm far}$, $\varepsilon_{\rm split}$, and interior/boundary residual norms.}
  \label{tab:tt-splitting-qc}
\end{table}

Table~\ref{tab:tt-splitting-qc} shows that both routes satisfy $q_{\rm far}(x_s)\approx P_{\rm far}$, with $\varepsilon_{\rm split}$ at $10^{-8}$ and $10^{-9}$ scales, while interior Bellman residuals stay at the $10^{-11}$ scale. This closes the mechanism loop behind the fast/slow peak split for both representative and showcase-quality clear-instance routes. We track these residuals together with $\varepsilon_{\rm split}$ as automatic QC scalars in continuous optimization runs.

\subsection*{Committor Implementation in This Report}
\begin{enumerate}
  \item Build the exact transition kernel for each representative case (including laziness, drift, membrane permeability, and reflecting boundaries).
  \item Use $A=\{x\le x_s+1,\ |y-y_{\rm mid}|\le 1\}$ and $B=\{\text{target}\}$, and solve $q$ via fixed-point iteration with tolerance $10^{-11}$; record interior Bellman residual and boundary mismatch as quality diagnostics.
  \item Use exact augmented-state recursion on $(x,\ell,s,r)$ (not MC): $\ell$ tracks leak, $s$ tracks whether $\Sigma_{q^*}$ has been visited, and $r$ tracks strict recross (return to $A$ after visiting $\Sigma_{q^*}$).
  \item Report $q^*=0.5$ as default and additionally provide $q^*\in\{0.4,0.5,0.6\}$ robustness results (Table~\ref{tab:membrane-qstar-sens}).
  \item For no-corridor two-target cases, keep the strict identity $f=f_{\rm near}+f_{\rm far}$ as level-1 decomposition, then interpret windows with committor-consistent labels.
\end{enumerate}

\section{Takeaways and Next Steps (Committor-Centric)}
\subsection*{Implementation Summary}
\begin{enumerate}
  \item Membrane corridor track: clear-double persists in both symmetric and asymmetric representatives; under strict recross, peak1 is mostly \texttt{L0R0}, valley gains \texttt{L1R0}, and peak2 is mostly \texttt{L0R1}.
  \item No-corridor two-target track: we now run ``all-candidates scan $\rightarrow$ adaptive local refinement $\rightarrow$ robust-gated clear-instance selection'' and report a dedicated figure/table pair (Figure~\ref{fig:tt-near-clear-instance}, Table~\ref{tab:tt-near-clear-instance}); the selected clear-instance is $(b_x,d,\Delta y)=(0.12,4,1)$, keeps positive phase=2 margins, and passes the first robustness gate ($\mathrm{peak\mbox{-}margin}\ge 0.07$, $\mathrm{min\mbox{-}margin}\ge 0.07$, $\mathrm{valley}/\max\le 0.10$). We additionally export strong-gate audit scalars (\texttt{gate-pass}/\texttt{gate-slack}) plus first-gate candidate count (168) and keep the hard clear-figure integrity check (file exists and size $>2$KB), with figure path/size metadata written to `summary.json` for continuous monitoring.
  \item Methodologically, committor is now fully operational in the pipeline; splitting committor is numerically closed against exact far-hit mass for both representative and clear-instance cases, and we jointly monitor $\varepsilon_{\rm split}$ plus committor residual diagnostics in automatic QC (current $\varepsilon_{\rm split}$ scale $\sim10^{-8}$ for representative and $\sim10^{-9}$ for clear-instance).
  \item Operationally, the current full local regression (codegen + CN/EN build checks) takes about $70\,$s per round (around 1.2 minutes), so it stays in the local-first regime ($<10$ minutes); we switch to Isambard only when projected single-round runtime exceeds $20$ minutes.
\end{enumerate}

\subsection*{Next Directions (All Grounded in Committor)}
\begin{enumerate}
  \item Scale exact decomposition from representative cases to full scan maps, and export threshold-margin maps for clear-instance quality (not just single-point tables).
  \item Run $q^*$ sensitivity analysis (e.g., $0.4/0.5/0.6$) and report peak/valley class stability.
  \item Systematically output splitting committor fields and window summaries for two-target scans:
  \[
  q_{\rm far}(x)=\mathbb P_x(\tau_{\rm far}<\tau_{\rm near}),
  \]
  and use both $\varepsilon_{\rm split}=|q_{\rm far}(x_s)-P_{\rm far}|$ (for representative and clear-instance routes) and \texttt{min-margin}/\texttt{gate-slack} as automatic quality checks.
  \item Compare top/bottom membrane asymmetry versus directional asymmetry (inside/outside permeability mismatch) in a unified committor geometry view.
  \item Keep continuous build checks active (codegen + EN/CN LaTeX) and enforce routing policy in keepalive: run local for short rounds ($<10$ minutes), and use isbard doctor/auth + submit/status/fetch for large scans (projected round $>20$ minutes), with automatic local-smoke fallback on remote failures.
  \item For every remote round, always record \texttt{JOB\_ID}, \texttt{REMOTE\_DIR}, key artifact paths, and explicit next actions in the run summary.
\end{enumerate}

\section*{Reproducibility}
\begin{verbatim}
cd reports/grid2d_membrane_near_target
../../.venv/bin/python code/membrane_near_target_report.py
latexmk -pdf -interaction=nonstopmode -halt-on-error \
  -auxdir=build -emulate-aux-dir grid2d_membrane_near_target_en.tex
\end{verbatim}

\end{document}
